\chapter{Referenced foundations and full working proofs}
\label{sec:proof-foundations}

The results below supply the mathematical foundations used in this report.
The explicitly assumed inputs are finite second moments and, where stated,
a jointly Gaussian law and independent observations. An assumption is not a
consequence of a simulation. Definitions, source quotations, historical
records, rounded numerical evaluations and recommendations have the status
specified in the claim map; they are not additional mathematical theorems.
The proofs reduce the deductions to linearity of integration, the spectral
theorem, the Gaussian characteristic function and conditional expectation.
Source locators refer to the linked author versions specified in the
bibliography, rather than silently mixing their pagination with print editions.

\noindent\textbf{Document-revision provenance.}
The current book-format revision is recorded in
\path{reporting/book-format/revision.json}, with current TeX/Bib sources
and PDF hashes. The earlier proof-text revision remains recorded in
\path{reporting/proof-revision/revision.json}.
The original \path{reporting/run.json} retains the numerical
simulation campaign, its original document hashes and numerical outputs.
The book revision checks and reuses the unchanged simulation-data files; it does
not rerun the experiment or claim a new native SPARK/Rocq proof.
The repository's subsequent shared-library refactor is separate from the
retained native campaigns. Their source-pinned results do not constitute
new GNATprove evidence for refactored Ada bodies; those changes require
their own matching campaign. The written mathematical proofs below do not
depend on that program-verification status.

\begin{proposition}[Finite moments and affine changes of scale]
\label{pf:moments}
Let a real random vector \(X\) have finite second moments, mean \(\mu\)
and covariance \(\Sigma\). Let \(Y\) be square-integrable, \(w\) a
deterministic vector of matching dimension, and \(b,s\) real constants. Then
\[
 \E(w^{\mathsf T}X)=w^{\mathsf T}\mu,\qquad
 \Var(w^{\mathsf T}X)=w^{\mathsf T}\Sigma w,\qquad
 \SD(b+sY)=|s|\SD(Y).
\]
A zero variance is equivalent to equality with the mean almost surely.
\end{proposition}
\begin{proof}
H\"older's inequality ensures integrability of all products
\((X_i-\mu_i)(X_j-\mu_j)\). Finite linearity of the integral then gives
the expectation formula and
\[
 \E\!\left[\Big(\sum_iw_i(X_i-\mu_i)\Big)^2\right]
 =\sum_{i,j}w_iw_j\E[(X_i-\mu_i)(X_j-\mu_j)]
 =w^{\mathsf T}\Sigma w.
\]
Replacing \(Y\) by \(b+sY\) multiplies its centered value by \(s\), so
its variance is \(s^2\Var(Y)\); taking the nonnegative square root proves
the SD formula. If \(\E[(Y-\E Y)^2]=0\), Markov's inequality applied to
\((Y-\E Y)^2\) makes \(P(|Y-\E Y|>\epsilon)=0\) for every
\(\epsilon>0\). The countable union over \(\epsilon=1/n\) proves the
almost-sure assertion. The reverse implication follows by integration.
These are applications of linearity, H\"older and the variance identities
in \citet[Section 1.6, especially Theorems 1.6.1, 1.6.3--1.6.4]{durrett2019}.
\end{proof}

\begin{proposition}[Covariance feasibility and Gaussian construction]
\label{pf:spectral}
A real symmetric matrix is a covariance matrix if and only if it is
positive semidefinite. For the equicorrelation matrix
\(R_n=(1-c)I_n+c\mathbf1_n\mathbf1_n^{\mathsf T}\), \(n\geq2\), this
is equivalent to \(-1/(n-1)\leq c\leq1\). Positive definiteness requires
both inequalities to be strict.
\end{proposition}
\begin{proof}
Proposition~\ref{pf:moments} gives
\(w^{\mathsf T}\Sigma w=\E[(w^{\mathsf T}(X-\mu))^2]\geq0\).
Conversely, the real spectral theorem diagonalizes a symmetric positive
semidefinite matrix as \(QDQ^{\mathsf T}\), with orthogonal \(Q\) and
nonnegative diagonal entries \(D\). For independent standard normal
coordinates \(Z\), the vector \(QD^{1/2}Z\) has covariance
\(QD^{1/2}I D^{1/2}Q^{\mathsf T}=\Sigma\) and is Gaussian, including
zero eigenvalues. The spectral theorem and positive square roots are
\citet[Theorems 7.29 and 7.38]{axler2024}; the Gaussian construction is
also \citet[Theorems 3.4.4--3.4.5]{gallager2013}.
For \(R_n\), the span of \(\mathbf1_n\) has eigenvalue
\(1+(n-1)c\), and its orthogonal complement has eigenvalue \(1-c\),
because \(\mathbf1_n^{\mathsf T}w=0\) there. Every vector decomposes
into these two orthogonal subspaces. Nonnegativity, or positivity, of
both eigenvalues gives precisely the stated inequalities. If one fixed
\(c\) must work for arbitrarily large \(n\), the lower bounds force
\(c\geq0\); the finite-dimensional bounds alone permit negative \(c\).
\end{proof}

\begin{proposition}[Gaussian closure and block independence]
\label{pf:gaussian}
An affine image of a jointly Gaussian vector is jointly Gaussian. Its law
is determined by its mean and covariance. Two blocks of a jointly Gaussian
vector are independent whenever their cross-covariance is zero, including
singular Gaussian vectors.
\end{proposition}
\begin{proof}
For a Gaussian vector with mean \(\mu\) and covariance \(\Sigma\), its
characteristic function is
\(\psi_X(t)=\exp(i t^{\mathsf T}\mu-t^{\mathsf T}\Sigma t/2)\).
Thus \(\psi_{AX+b}(t)=e^{it^{\mathsf T}b}\psi_X(A^{\mathsf T}t)\)
has the Gaussian form with mean \(A\mu+b\) and covariance
\(A\Sigma A^{\mathsf T}\). This also follows from the independent-normal
affine representation, without invertibility. If a two-block covariance
has zero off-diagonal blocks, the displayed characteristic function
factors as \(\psi_U(s)\psi_V(t)\). By Fubini this is the characteristic
function of the product of the two marginal measures. Uniqueness of
characteristic functions proves both independence and determination by
mean and covariance. The Gaussian formula and closure are
\citet[Theorems 3.3.1--3.3.2, equations (3.14) and (3.18)]{gallager2013};
uniqueness is \citet[Theorem 3.10.4]{durrett2019}.
\end{proof}

\begin{proposition}[Gaussian conditional distribution, including exact histories]
\label{pf:conditioning}
Suppose \((U,V)\) is jointly Gaussian with mean \((\mu_U,\mu_V)\),
covariance blocks \(A,B\) and \(C=\Cov(U,V)\), and \(B\) positive
definite. A version of the conditional law, for every \(v\), is Gaussian
with mean and covariance
\[
 \mu_U+CB^{-1}(v-\mu_V),\qquad A-CB^{-1}C^{\mathsf T}.
\]
The conditional-mean formula is the unique continuous version on the
whole conditioning space. The full joint covariance need not be invertible.
\end{proposition}
\begin{proof}
Set \(L=CB^{-1}\) and
\(W=U-\mu_U-L(V-\mu_V)\). Its cross-covariance with \(V\) is
\(C-LB=0\). Expanding its covariance gives
\(A-LC^{\mathsf T}-CL^{\mathsf T}+LBL^{\mathsf T}
=A-CB^{-1}C^{\mathsf T}\). Proposition~\ref{pf:gaussian} makes
\(W\) Gaussian and independent of \(V\). Consequently the kernel
\(K(v,H)=P(\mu_U+L(v-\mu_V)+W\in H)\) satisfies
\[
 P(U\in H,V\in D)=\int_D K(v,H)P_V(dv)
\]
by independence and Fubini, which is the defining conditional-law identity.
Its expectation is the displayed affine function. A positive-definite
Gaussian \(V\) has a strictly positive density everywhere. If two
continuous versions differed at any point, continuity would give a
neighborhood of positive probability on which they differ, contradicting
almost-sure uniqueness. Conditioning on an exact continuous history uses
this specified version, not division by the probability-zero event
\(\{V=v\}\). This residual proof is the Gaussian regression theorem
\citep[Theorems 3.5.1--3.5.2]{gallager2013}, proved here in the slightly
more general form allowing singular \((U,V)\); conditional versions are
defined in \citet[Section 4.1]{durrett2019}.
\end{proof}

\begin{proposition}[Normal marginal distributions do not imply joint normality]
\label{pf:normal-marginals}
Normal marginals and a specified Pearson correlation do not in general
determine Gaussian composite or joint-tail laws.
\end{proposition}
\begin{proof}
Let \(X\sim N(0,1)\) and independently let \(S\) equal \(1\) with
probability \((1+\rho)/2\) and \(-1\) otherwise, \(-1<\rho<1\).
Put \(Y=SX\). Symmetry gives \(Y\sim N(0,1)\), and independence gives
\(\E[XY]=\E[S]\E[X^2]=\rho\). Nevertheless
\(P(X+Y=0)=(1-\rho)/2>0\). A nondegenerate normal variable has a
density and no atoms; \(X+Y\) has variance \(2+2\rho>0\), so this
pair is not jointly Gaussian. For a concrete tail nonidentification at
\(\rho=0\), let \(t>0\) and \(p=P(X>t)\in(0,1/2)\). This pair
has \(P(X>t,Y>t)=p/2\), whereas an independent standard-normal pair,
with the same marginals and zero correlation, has intersection \(p^2\).
These differ since \(p<1/2\). The construction uses symmetry,
independence and the product-integration rule
\citep[Sections 1.7 and 2.1]{durrett2019}; it is a counterexample,
not an alternative assumption secretly used for the report's Gaussian tables.
\end{proof}

\begin{proposition}[Event arithmetic and conditional denominators]
\label{pf:events}
For events \(A,B\) with probabilities \(p_A,p_B\) and intersection
\(J\), their union has probability \(p_A+p_B-J\). If its conditioning
probability is positive, \(P(B\mid A)=J/p_A\). The correlation of the
nonconstant indicators is
\[
 \frac{J-p_Ap_B}{\sqrt{p_A(1-p_A)p_B(1-p_B)}}.
\]
\end{proposition}
\begin{proof}
The pointwise identities
\(1_{A\cup B}=1_A+1_B-1_{A\cap B}\) and
\(1_A1_B=1_{A\cap B}\) prove the union and covariance formulas by
linearity. Because \(1_A^2=1_A\),
\(\Var(1_A)=p_A-p_A^2\), and similarly for \(B\); divide their
covariance by the product of their positive SDs. The conditional formula
is precisely the event-conditional probability definition, requiring
\(p_A>0\). The four disjoint cells have probabilities
\(J,p_A-J,p_B-J,1-p_A-p_B+J\), obtained by the same identities.
These are consequences of finite additivity and conditional probability
\citep[Sections 1.1 and 4.1]{durrett2019}.
\end{proof}

\begin{proposition}[Tower property and population--conditional variance]
\label{pf:expectation}
For square-integrable \(Y\) and a sub-\(\sigma\)-algebra \(\mathcal G\),
\[
 \E\E[Y\mid\mathcal G]=\E Y,\qquad
 \Var(Y)=\E\Var(Y\mid\mathcal G)
              +\Var(\E[Y\mid\mathcal G]).
\]
\end{proposition}
\begin{proof}
The defining conditional-integral identity with the whole space gives
the first formula. Write \(M=\E[Y\mid\mathcal G]\). Conditional
Jensen \citep[Theorem 4.1.10]{durrett2019} ensures \(M\) is
square-integrable. In
\(Y-\E Y=(Y-M)+(M-\E Y)\), the cross term has expectation zero:
the pull-out and tower properties give
\(\E[(Y-M)(M-\E Y)]
=\E[(M-\E Y)\E(Y-M\mid\mathcal G)]=0\).
Expanding the square and integrating gives the variance identity. The
pull-out rule extends from bounded multipliers by truncation and
Cauchy--Schwarz. These conditional-expectation properties follow from
\citet[Section 4.1, especially Theorems 4.1.13--4.1.14]{durrett2019}.
In particular, fixing an observed starting score and varying it across
the population are different probability experiments, even when both
use the same marginal test SD.
\end{proof}
